\documentclass[11pt,a4paper]{article}
% lesson-preamble.tex -- Shared preamble for all Phase 00 lessons
% Usage: % lesson-preamble.tex -- Shared preamble for all Phase 00 lessons
% Usage: % lesson-preamble.tex -- Shared preamble for all Phase 00 lessons
% Usage: \input{../lesson-preamble} (from subdomain directories)

% ---- Encoding and fonts ----
\usepackage[utf8]{inputenc}
\usepackage[T1]{fontenc}

% ---- Mathematics ----
\usepackage{amsmath,amssymb,amsthm}
\usepackage{mathtools}

% ---- Page geometry ----
\usepackage[margin=1in,headheight=14pt]{geometry}

% ---- Hyperlinks ----
\usepackage[colorlinks=true,linkcolor=red,citecolor=blue,urlcolor=blue]{hyperref}

% ---- Lists ----
\usepackage{enumitem}

% ---- Colors and boxes ----
\usepackage{xcolor}
\usepackage[theorems,skins,breakable]{tcolorbox}

% ---- Header/Footer ----
\usepackage{fancyhdr}
\renewcommand{\headrulewidth}{0.4pt}

% ---- Tables ----
\usepackage{booktabs}

% ---- Bibliography ----
\usepackage{natbib}
\bibliographystyle{plainnat}

% --------------------------------------------------------------------------
% Theorem environments
% --------------------------------------------------------------------------
\theoremstyle{plain}
\newtheorem{theorem}{Theorem}[section]
\newtheorem{lemma}[theorem]{Lemma}
\newtheorem{proposition}[theorem]{Proposition}
\newtheorem{corollary}[theorem]{Corollary}

\theoremstyle{definition}
\newtheorem{definition}[theorem]{Definition}
\newtheorem{example}[theorem]{Example}
\newtheorem{exercise}[theorem]{Exercise}
\newtheorem{assumption}[theorem]{Assumption}

\theoremstyle{remark}
\newtheorem{remark}[theorem]{Remark}

% --------------------------------------------------------------------------
% Colored boxes
% --------------------------------------------------------------------------
\newtcolorbox{keyresult}[1][]{%
  colback=green!5!white,
  colframe=green!50!black,
  fonttitle=\bfseries,
  title={Key Result},
  breakable,
  #1
}

\newtcolorbox{rlconnection}[1][]{%
  colback=blue!5!white,
  colframe=blue!50!black,
  fonttitle=\bfseries,
  title={RL/ML Connection},
  breakable,
  #1
}

\newtcolorbox{warning}[1][]{%
  colback=red!5!white,
  colframe=red!50!black,
  fonttitle=\bfseries,
  title={Warning},
  breakable,
  #1
}

% --------------------------------------------------------------------------
% Custom commands -- number sets
% --------------------------------------------------------------------------
\newcommand{\R}{\mathbb{R}}
\newcommand{\N}{\mathbb{N}}
\newcommand{\Q}{\mathbb{Q}}
\newcommand{\Z}{\mathbb{Z}}
\newcommand{\C}{\mathbb{C}}
\newcommand{\F}{\mathbb{F}}

% --------------------------------------------------------------------------
% Custom commands -- probability and expectation
% --------------------------------------------------------------------------
\newcommand{\E}{\mathbb{E}}
\newcommand{\Prob}{\mathbb{P}}
\newcommand{\Var}{\operatorname{Var}}
\newcommand{\Cov}{\operatorname{Cov}}
\newcommand{\indicator}{\mathbf{1}}

% --------------------------------------------------------------------------
% Custom commands -- convergence arrows
% --------------------------------------------------------------------------
\newcommand{\cas}{\xrightarrow{\text{a.s.}}}
\newcommand{\cprob}{\xrightarrow{\Prob}}
\newcommand{\clp}[1]{\xrightarrow{L^{#1}}}
\newcommand{\cdist}{\xrightarrow{d}}

% --------------------------------------------------------------------------
% Custom commands -- common operators
% --------------------------------------------------------------------------
\DeclareMathOperator{\dom}{dom}
\DeclareMathOperator{\epi}{epi}
\DeclareMathOperator{\conv}{conv}
\DeclareMathOperator{\relint}{relint}
\DeclareMathOperator{\aff}{aff}
\DeclareMathOperator{\cl}{cl}
\DeclareMathOperator{\interior}{int}
\DeclareMathOperator{\tr}{tr}
\DeclareMathOperator{\diag}{diag}
\DeclareMathOperator{\rank}{rank}
\DeclareMathOperator{\sgn}{sgn}
\DeclareMathOperator{\supp}{supp}
\DeclareMathOperator{\argmin}{arg\,min}
\DeclareMathOperator{\argmax}{arg\,max}
\DeclareMathOperator{\prox}{prox}
\DeclareMathOperator{\dist}{dist}
\DeclareMathOperator{\diam}{diam}
\DeclareMathOperator{\Span}{span}

% --------------------------------------------------------------------------
% Custom commands -- linear algebra operators
% --------------------------------------------------------------------------
\DeclareMathOperator{\nullity}{nullity}
\DeclareMathOperator{\range}{range}
\DeclareMathOperator{\nullspace}{null}
\DeclareMathOperator{\proj}{proj}

% --------------------------------------------------------------------------
% Custom commands -- measure theory
% --------------------------------------------------------------------------
\newcommand{\powerset}{\mathcal{P}}
\newcommand{\Borel}{\mathcal{B}}
 (from subdomain directories)

% ---- Encoding and fonts ----
\usepackage[utf8]{inputenc}
\usepackage[T1]{fontenc}

% ---- Mathematics ----
\usepackage{amsmath,amssymb,amsthm}
\usepackage{mathtools}

% ---- Page geometry ----
\usepackage[margin=1in,headheight=14pt]{geometry}

% ---- Hyperlinks ----
\usepackage[colorlinks=true,linkcolor=red,citecolor=blue,urlcolor=blue]{hyperref}

% ---- Lists ----
\usepackage{enumitem}

% ---- Colors and boxes ----
\usepackage{xcolor}
\usepackage[theorems,skins,breakable]{tcolorbox}

% ---- Header/Footer ----
\usepackage{fancyhdr}
\renewcommand{\headrulewidth}{0.4pt}

% ---- Tables ----
\usepackage{booktabs}

% ---- Bibliography ----
\usepackage{natbib}
\bibliographystyle{plainnat}

% --------------------------------------------------------------------------
% Theorem environments
% --------------------------------------------------------------------------
\theoremstyle{plain}
\newtheorem{theorem}{Theorem}[section]
\newtheorem{lemma}[theorem]{Lemma}
\newtheorem{proposition}[theorem]{Proposition}
\newtheorem{corollary}[theorem]{Corollary}

\theoremstyle{definition}
\newtheorem{definition}[theorem]{Definition}
\newtheorem{example}[theorem]{Example}
\newtheorem{exercise}[theorem]{Exercise}
\newtheorem{assumption}[theorem]{Assumption}

\theoremstyle{remark}
\newtheorem{remark}[theorem]{Remark}

% --------------------------------------------------------------------------
% Colored boxes
% --------------------------------------------------------------------------
\newtcolorbox{keyresult}[1][]{%
  colback=green!5!white,
  colframe=green!50!black,
  fonttitle=\bfseries,
  title={Key Result},
  breakable,
  #1
}

\newtcolorbox{rlconnection}[1][]{%
  colback=blue!5!white,
  colframe=blue!50!black,
  fonttitle=\bfseries,
  title={RL/ML Connection},
  breakable,
  #1
}

\newtcolorbox{warning}[1][]{%
  colback=red!5!white,
  colframe=red!50!black,
  fonttitle=\bfseries,
  title={Warning},
  breakable,
  #1
}

% --------------------------------------------------------------------------
% Custom commands -- number sets
% --------------------------------------------------------------------------
\newcommand{\R}{\mathbb{R}}
\newcommand{\N}{\mathbb{N}}
\newcommand{\Q}{\mathbb{Q}}
\newcommand{\Z}{\mathbb{Z}}
\newcommand{\C}{\mathbb{C}}
\newcommand{\F}{\mathbb{F}}

% --------------------------------------------------------------------------
% Custom commands -- probability and expectation
% --------------------------------------------------------------------------
\newcommand{\E}{\mathbb{E}}
\newcommand{\Prob}{\mathbb{P}}
\newcommand{\Var}{\operatorname{Var}}
\newcommand{\Cov}{\operatorname{Cov}}
\newcommand{\indicator}{\mathbf{1}}

% --------------------------------------------------------------------------
% Custom commands -- convergence arrows
% --------------------------------------------------------------------------
\newcommand{\cas}{\xrightarrow{\text{a.s.}}}
\newcommand{\cprob}{\xrightarrow{\Prob}}
\newcommand{\clp}[1]{\xrightarrow{L^{#1}}}
\newcommand{\cdist}{\xrightarrow{d}}

% --------------------------------------------------------------------------
% Custom commands -- common operators
% --------------------------------------------------------------------------
\DeclareMathOperator{\dom}{dom}
\DeclareMathOperator{\epi}{epi}
\DeclareMathOperator{\conv}{conv}
\DeclareMathOperator{\relint}{relint}
\DeclareMathOperator{\aff}{aff}
\DeclareMathOperator{\cl}{cl}
\DeclareMathOperator{\interior}{int}
\DeclareMathOperator{\tr}{tr}
\DeclareMathOperator{\diag}{diag}
\DeclareMathOperator{\rank}{rank}
\DeclareMathOperator{\sgn}{sgn}
\DeclareMathOperator{\supp}{supp}
\DeclareMathOperator{\argmin}{arg\,min}
\DeclareMathOperator{\argmax}{arg\,max}
\DeclareMathOperator{\prox}{prox}
\DeclareMathOperator{\dist}{dist}
\DeclareMathOperator{\diam}{diam}
\DeclareMathOperator{\Span}{span}

% --------------------------------------------------------------------------
% Custom commands -- linear algebra operators
% --------------------------------------------------------------------------
\DeclareMathOperator{\nullity}{nullity}
\DeclareMathOperator{\range}{range}
\DeclareMathOperator{\nullspace}{null}
\DeclareMathOperator{\proj}{proj}

% --------------------------------------------------------------------------
% Custom commands -- measure theory
% --------------------------------------------------------------------------
\newcommand{\powerset}{\mathcal{P}}
\newcommand{\Borel}{\mathcal{B}}
 (from subdomain directories)

% ---- Encoding and fonts ----
\usepackage[utf8]{inputenc}
\usepackage[T1]{fontenc}

% ---- Mathematics ----
\usepackage{amsmath,amssymb,amsthm}
\usepackage{mathtools}

% ---- Page geometry ----
\usepackage[margin=1in,headheight=14pt]{geometry}

% ---- Hyperlinks ----
\usepackage[colorlinks=true,linkcolor=red,citecolor=blue,urlcolor=blue]{hyperref}

% ---- Lists ----
\usepackage{enumitem}

% ---- Colors and boxes ----
\usepackage{xcolor}
\usepackage[theorems,skins,breakable]{tcolorbox}

% ---- Header/Footer ----
\usepackage{fancyhdr}
\renewcommand{\headrulewidth}{0.4pt}

% ---- Tables ----
\usepackage{booktabs}

% ---- Bibliography ----
\usepackage{natbib}
\bibliographystyle{plainnat}

% --------------------------------------------------------------------------
% Theorem environments
% --------------------------------------------------------------------------
\theoremstyle{plain}
\newtheorem{theorem}{Theorem}[section]
\newtheorem{lemma}[theorem]{Lemma}
\newtheorem{proposition}[theorem]{Proposition}
\newtheorem{corollary}[theorem]{Corollary}

\theoremstyle{definition}
\newtheorem{definition}[theorem]{Definition}
\newtheorem{example}[theorem]{Example}
\newtheorem{exercise}[theorem]{Exercise}
\newtheorem{assumption}[theorem]{Assumption}

\theoremstyle{remark}
\newtheorem{remark}[theorem]{Remark}

% --------------------------------------------------------------------------
% Colored boxes
% --------------------------------------------------------------------------
\newtcolorbox{keyresult}[1][]{%
  colback=green!5!white,
  colframe=green!50!black,
  fonttitle=\bfseries,
  title={Key Result},
  breakable,
  #1
}

\newtcolorbox{rlconnection}[1][]{%
  colback=blue!5!white,
  colframe=blue!50!black,
  fonttitle=\bfseries,
  title={RL/ML Connection},
  breakable,
  #1
}

\newtcolorbox{warning}[1][]{%
  colback=red!5!white,
  colframe=red!50!black,
  fonttitle=\bfseries,
  title={Warning},
  breakable,
  #1
}

% --------------------------------------------------------------------------
% Custom commands -- number sets
% --------------------------------------------------------------------------
\newcommand{\R}{\mathbb{R}}
\newcommand{\N}{\mathbb{N}}
\newcommand{\Q}{\mathbb{Q}}
\newcommand{\Z}{\mathbb{Z}}
\newcommand{\C}{\mathbb{C}}
\newcommand{\F}{\mathbb{F}}

% --------------------------------------------------------------------------
% Custom commands -- probability and expectation
% --------------------------------------------------------------------------
\newcommand{\E}{\mathbb{E}}
\newcommand{\Prob}{\mathbb{P}}
\newcommand{\Var}{\operatorname{Var}}
\newcommand{\Cov}{\operatorname{Cov}}
\newcommand{\indicator}{\mathbf{1}}

% --------------------------------------------------------------------------
% Custom commands -- convergence arrows
% --------------------------------------------------------------------------
\newcommand{\cas}{\xrightarrow{\text{a.s.}}}
\newcommand{\cprob}{\xrightarrow{\Prob}}
\newcommand{\clp}[1]{\xrightarrow{L^{#1}}}
\newcommand{\cdist}{\xrightarrow{d}}

% --------------------------------------------------------------------------
% Custom commands -- common operators
% --------------------------------------------------------------------------
\DeclareMathOperator{\dom}{dom}
\DeclareMathOperator{\epi}{epi}
\DeclareMathOperator{\conv}{conv}
\DeclareMathOperator{\relint}{relint}
\DeclareMathOperator{\aff}{aff}
\DeclareMathOperator{\cl}{cl}
\DeclareMathOperator{\interior}{int}
\DeclareMathOperator{\tr}{tr}
\DeclareMathOperator{\diag}{diag}
\DeclareMathOperator{\rank}{rank}
\DeclareMathOperator{\sgn}{sgn}
\DeclareMathOperator{\supp}{supp}
\DeclareMathOperator{\argmin}{arg\,min}
\DeclareMathOperator{\argmax}{arg\,max}
\DeclareMathOperator{\prox}{prox}
\DeclareMathOperator{\dist}{dist}
\DeclareMathOperator{\diam}{diam}
\DeclareMathOperator{\Span}{span}

% --------------------------------------------------------------------------
% Custom commands -- linear algebra operators
% --------------------------------------------------------------------------
\DeclareMathOperator{\nullity}{nullity}
\DeclareMathOperator{\range}{range}
\DeclareMathOperator{\nullspace}{null}
\DeclareMathOperator{\proj}{proj}

% --------------------------------------------------------------------------
% Custom commands -- measure theory
% --------------------------------------------------------------------------
\newcommand{\powerset}{\mathcal{P}}
\newcommand{\Borel}{\mathcal{B}}

\pagestyle{fancy}
\fancyhf{}
\fancyhead[L]{\textsc{Lesson 3: Differentiation and Integration}}
\fancyhead[R]{\thepage}
\title{Lesson 3: Differentiation and Integration}
\author{AI/ML Mathematics}
\date{}
\begin{document}
\maketitle
\tableofcontents
\newpage

%% ============================================================
%% SECTION 1: MOTIVATION AND OVERVIEW
%% ============================================================
\section{Motivation and Overview}
\label{sec:motivation}

Every modern deep learning system is trained by \emph{gradient descent}:
given a loss function $L(\theta)$, we compute its derivative with respect
to the parameters $\theta$ and update $\theta \leftarrow \theta -
\eta \nabla L(\theta)$. The chain rule, applied layer by layer through
the network, is the mathematical engine behind backpropagation. But what
guarantees that the derivative exists, that the chain rule is valid, and
that the loss decrease can be quantified? The \emph{mean value theorem}
gives precise bounds on how much a function changes in terms of its
derivative, and \emph{Taylor's theorem} provides the polynomial
approximation that justifies gradient descent convergence rates.

On the integration side, expected values $\E[f(X)] = \int f(x)\, p(x)
\,dx$ are the foundation of every statistical learning objective, from
maximum likelihood to policy gradient. The Riemann integral provides the
first rigorous framework for these integrals, and the \emph{fundamental
theorem of calculus} connects differentiation and integration, enabling
the computation of integrals via antiderivatives. The tools developed in
this lesson will let us make all of these arguments precise.

\begin{rlconnection}[title={Why RL/ML needs differentiation and integration}]
\begin{itemize}[nosep]
  \item \textbf{Backpropagation.} The chain rule computes
    $\partial L / \partial \theta_i$ by composing derivatives through
    layers of a neural network, making gradient-based training possible.
  \item \textbf{Convergence of gradient descent.} The mean value theorem
    and Taylor's theorem establish that gradient descent with step size
    $\eta$ decreases the loss by approximately
    $\eta \|\nabla L\|^2$, yielding convergence rates.
  \item \textbf{Policy gradient.} The REINFORCE estimator
    $\nabla_\theta J(\theta) = \E[\nabla_\theta \log \pi_\theta(a|s)
    \cdot Q^\pi(s,a)]$ requires exchanging differentiation and
    integration, justified by the dominated convergence theorem (which
    extends the Riemann integral framework developed here).
  \item \textbf{Limitations of Riemann integration.} Many functions
    arising in probability (e.g., indicator functions of complicated
    sets) are not Riemann integrable, motivating the Lebesgue integral
    in Phase~01.
\end{itemize}
\end{rlconnection}

\paragraph{Prerequisites.}
Continuity ($\varepsilon$-$\delta$ and sequential characterizations), the
extreme value theorem, the intermediate value theorem, and uniform
continuity (Lesson~2). Completeness of $\R$, convergence of sequences,
and supremum/infimum (Lesson~1).

\paragraph{Learning objectives.}
After completing this lesson, you will be able to:
\begin{itemize}[nosep]
  \item \textbf{Define} the derivative and \textbf{prove} that
    differentiability implies continuity.
  \item \textbf{Derive} the sum, product, quotient, and chain rules for
    differentiation.
  \item \textbf{State} and \textbf{prove} Fermat's theorem, Rolle's
    theorem, and the mean value theorem.
  \item \textbf{Apply} Taylor's theorem with Lagrange remainder to obtain
    polynomial approximations and error bounds.
  \item \textbf{Define} the Riemann integral via partitions and upper/lower
    sums, and \textbf{prove} that continuous functions are integrable.
  \item \textbf{State} and \textbf{prove} the fundamental theorem of
    calculus (both parts).
\end{itemize}

%% ============================================================
%% SECTION 2: REFERENCES
%% ============================================================
\section{References}
\label{sec:references}

\begin{itemize}
  \item \citet[Chapters 5--6]{rudin1976} --- Rigorous treatment of
    differentiation and the Riemann--Stieltjes integral.
  \item \citet[Chapters 5--7]{abbott2015} --- Accessible development of
    derivatives, the mean value theorem, and Riemann integration.
  \item \citet[Chapters 6--7]{bartle2011} --- Detailed coverage of
    differentiation rules, Taylor's theorem, and the Riemann integral
    with many examples.
  \item \citet[Chapters 5--6]{pugh2015} --- Differentiation and
    integration with a geometric viewpoint.
  \item \citet[Chapter 1]{folland1999} --- Brief review of Riemann
    integration as motivation for the Lebesgue theory.
\end{itemize}

%% ============================================================
%% SECTION 3: THE DERIVATIVE
%% ============================================================
\section{The Derivative}
\label{sec:derivative}

\begin{definition}[Derivative]
\label{def:derivative}
Let $f \colon (a, b) \to \R$ and $c \in (a, b)$. The \emph{derivative}
of $f$ at $c$ is
\[
  f'(c) = \lim_{x \to c} \frac{f(x) - f(c)}{x - c},
\]
provided this limit exists. If $f'(c)$ exists, we say $f$ is
\emph{differentiable} at $c$. Equivalently, $f'(c) = \lim_{h \to 0}
(f(c + h) - f(c))/h$.
\end{definition}

\begin{theorem}[Differentiability Implies Continuity]
\label{thm:diff_implies_cont}
If $f$ is differentiable at $c$, then $f$ is continuous at $c$.
\end{theorem}

\begin{proof}
Suppose $f'(c)$ exists. For $x \neq c$:
\[
  f(x) - f(c) = \frac{f(x) - f(c)}{x - c} \cdot (x - c).
\]
Taking $x \to c$: the first factor converges to $f'(c)$ and the second
to $0$, so $f(x) - f(c) \to 0$ by the algebraic limit theorems
(Lesson~1, Theorem~4.4). Hence $f(x) \to f(c)$, i.e., $f$ is continuous
at $c$.
\end{proof}

\begin{example}
\label{ex:abs_not_diff}
The function $f(x) = |x|$ is continuous at $0$ but not differentiable
there. The left-hand limit $\lim_{h \to 0^-} (|h| - 0)/h = -1$ and the
right-hand limit $\lim_{h \to 0^+} |h|/h = 1$ disagree, so the
two-sided limit does not exist.
\end{example}

\begin{example}
\label{ex:power_rule}
For $f(x) = x^n$ with $n \in \N$, the derivative is $f'(x) = nx^{n-1}$.
Indeed,
\[
  \frac{(c+h)^n - c^n}{h}
  = \frac{h\bigl(nc^{n-1} + \binom{n}{2}c^{n-2}h + \cdots + h^{n-1}\bigr)}{h}
  \to nc^{n-1}
\]
as $h \to 0$.
\end{example}

%% ============================================================
%% SECTION 4: DIFFERENTIATION RULES
%% ============================================================
\section{Differentiation Rules}
\label{sec:diff_rules}

\begin{theorem}[Sum and Constant Multiple Rules]
\label{thm:sum_rule}
If $f$ and $g$ are differentiable at $c$ and $\alpha \in \R$, then:
\begin{enumerate}[nosep,label=(\roman*)]
  \item $(f + g)'(c) = f'(c) + g'(c)$.
  \item $(\alpha f)'(c) = \alpha f'(c)$.
\end{enumerate}
\end{theorem}

\begin{proof}
(i) $\displaystyle\frac{(f+g)(x) - (f+g)(c)}{x - c}
= \frac{f(x) - f(c)}{x - c} + \frac{g(x) - g(c)}{x - c}
\to f'(c) + g'(c)$ as $x \to c$, by the algebraic limit theorem for
sums.

(ii) is similar, using the constant multiple rule for limits.
\end{proof}

\begin{theorem}[Product Rule]
\label{thm:product_rule}
If $f$ and $g$ are differentiable at $c$, then
$(fg)'(c) = f'(c)\,g(c) + f(c)\,g'(c)$.
\end{theorem}

\begin{proof}
We write:
\begin{align*}
  \frac{f(x)g(x) - f(c)g(c)}{x - c}
  &= \frac{f(x)g(x) - f(c)g(x) + f(c)g(x) - f(c)g(c)}{x - c} \\
  &= \frac{f(x) - f(c)}{x - c}\,g(x) + f(c)\,\frac{g(x) - g(c)}{x - c}.
\end{align*}
As $x \to c$: the first term converges to $f'(c) \cdot g(c)$ (using
$g(x) \to g(c)$ by Theorem~\ref{thm:diff_implies_cont}), and the second
to $f(c) \cdot g'(c)$.
\end{proof}

\begin{theorem}[Quotient Rule]
\label{thm:quotient_rule}
If $f$ and $g$ are differentiable at $c$ and $g(c) \neq 0$, then
\[
  \left(\frac{f}{g}\right)'(c)
  = \frac{f'(c)\,g(c) - f(c)\,g'(c)}{g(c)^2}.
\]
\end{theorem}

\begin{proof}
It suffices to show $(1/g)'(c) = -g'(c)/g(c)^2$ and then apply the
product rule to $f \cdot (1/g)$.  Since $g(c) \neq 0$ and $g$ is
continuous at $c$, $g(x) \neq 0$ near $c$.  Then:
\[
  \frac{1/(g(x)) - 1/(g(c))}{x - c}
  = \frac{g(c) - g(x)}{(x - c)\,g(x)\,g(c)}
  = -\frac{1}{g(x)\,g(c)} \cdot \frac{g(x) - g(c)}{x - c}
  \to -\frac{g'(c)}{g(c)^2}. \qedhere
\]
\end{proof}

\begin{theorem}[Chain Rule]
\label{thm:chain_rule}
Let $f$ be differentiable at $c$ and $g$ be differentiable at $f(c)$.
Then $g \circ f$ is differentiable at $c$ and
\[
  (g \circ f)'(c) = g'(f(c)) \cdot f'(c).
\]
\end{theorem}

\begin{proof}
Define the auxiliary function
\[
  \phi(y) = \begin{cases}
    \dfrac{g(y) - g(f(c))}{y - f(c)} & \text{if } y \neq f(c), \\[6pt]
    g'(f(c)) & \text{if } y = f(c).
  \end{cases}
\]
Since $g$ is differentiable at $f(c)$, $\phi$ is continuous at $f(c)$.
For $x \neq c$ with $f(x) \neq f(c)$:
\[
  \frac{g(f(x)) - g(f(c))}{x - c}
  = \phi(f(x)) \cdot \frac{f(x) - f(c)}{x - c}.
\]
This identity also holds when $f(x) = f(c)$ (both sides are $0$). Taking
$x \to c$: $f(x) \to f(c)$ (continuity from differentiability), so
$\phi(f(x)) \to \phi(f(c)) = g'(f(c))$, and
$(f(x) - f(c))/(x - c) \to f'(c)$. Hence
$(g \circ f)'(c) = g'(f(c)) \cdot f'(c)$.
\end{proof}

\begin{rlconnection}[title={The chain rule as the backbone of backpropagation}]
A neural network computes a composition $f = f_L \circ f_{L-1} \circ
\cdots \circ f_1$, where each $f_\ell(\mathbf{x}) = \sigma(W_\ell
\mathbf{x} + \mathbf{b}_\ell)$ applies an affine transformation followed
by a nonlinearity $\sigma$. The chain rule gives
\[
  \frac{\partial L}{\partial W_\ell}
  = \frac{\partial L}{\partial f_L} \cdot
    \frac{\partial f_L}{\partial f_{L-1}} \cdots
    \frac{\partial f_{\ell+1}}{\partial f_\ell} \cdot
    \frac{\partial f_\ell}{\partial W_\ell},
\]
which is computed efficiently by the backpropagation algorithm.  Without
the chain rule, computing gradients of deep networks would require
finite-difference approximations with cost proportional to the number of
parameters.
\end{rlconnection}

\begin{example}
\label{ex:chain_rule}
Let $h(x) = (3x^2 + 1)^5$. Setting $f(x) = 3x^2 + 1$ and $g(y) = y^5$,
the chain rule gives $h'(x) = 5(3x^2+1)^4 \cdot 6x = 30x(3x^2+1)^4$.
\end{example}

%% ============================================================
%% SECTION 5: MEAN VALUE THEOREM AND CONSEQUENCES
%% ============================================================
\section{Mean Value Theorem and Consequences}
\label{sec:mvt}

\begin{theorem}[Fermat's Theorem]
\label{thm:fermat}
Let $f \colon (a, b) \to \R$ and suppose $f$ attains a local maximum or
minimum at $c \in (a, b)$. If $f$ is differentiable at $c$, then
$f'(c) = 0$.
\end{theorem}

\begin{proof}
Suppose $f$ has a local maximum at $c$: there exists $\delta > 0$ such
that $f(x) \le f(c)$ for $|x - c| < \delta$. For $0 < h < \delta$:
\[
  \frac{f(c + h) - f(c)}{h} \le 0,
\]
so $f'(c) = \lim_{h \to 0^+} (f(c+h) - f(c))/h \le 0$. For
$-\delta < h < 0$:
\[
  \frac{f(c + h) - f(c)}{h} \ge 0,
\]
so $f'(c) = \lim_{h \to 0^-} (f(c+h) - f(c))/h \ge 0$. Together,
$f'(c) = 0$. The proof for a local minimum is analogous.
\end{proof}

\begin{theorem}[Rolle's Theorem]
\label{thm:rolle}
If $f \colon [a, b] \to \R$ is continuous on $[a, b]$, differentiable on
$(a, b)$, and $f(a) = f(b)$, then there exists $c \in (a, b)$ with
$f'(c) = 0$.
\end{theorem}

\begin{proof}
If $f$ is constant on $[a, b]$, then $f'(c) = 0$ for every
$c \in (a, b)$. Otherwise, $f$ is not constant, so by the extreme value
theorem (Lesson~2, Theorem~6.2), $f$ attains its maximum $M$ and minimum
$m$ on $[a, b]$. Since $f$ is not constant and $f(a) = f(b)$, at least
one extreme value is attained at some $c \in (a, b)$ (both cannot occur
only at the endpoints, since $f(a) = f(b)$). By Fermat's theorem,
$f'(c) = 0$.
\end{proof}

\begin{keyresult}[title={Mean Value Theorem}]
\begin{theorem}[Mean Value Theorem (MVT)]
\label{thm:mvt}
If $f \colon [a, b] \to \R$ is continuous on $[a, b]$ and differentiable
on $(a, b)$, then there exists $c \in (a, b)$ with
\[
  f'(c) = \frac{f(b) - f(a)}{b - a}.
\]
\end{theorem}
\end{keyresult}

\begin{proof}
Define the auxiliary function
\[
  g(x) = f(x) - f(a) - \frac{f(b) - f(a)}{b - a}(x - a).
\]
Then $g$ is continuous on $[a, b]$, differentiable on $(a, b)$, and
$g(a) = 0 = g(b)$. By Rolle's theorem, there exists $c \in (a, b)$
with $g'(c) = 0$, i.e.,
\[
  f'(c) - \frac{f(b) - f(a)}{b - a} = 0. \qedhere
\]
\end{proof}

\begin{example}
\label{ex:mvt_speed}
If a car travels 150 km in 2 hours, the MVT guarantees there is an
instant when the car's speed is exactly $150/2 = 75$ km/h, regardless
of how the speed varied during the trip.
\end{example}

\begin{corollary}[Monotonicity from the Derivative]
\label{cor:monotone_deriv}
Let $f$ be continuous on $[a, b]$ and differentiable on $(a, b)$.
\begin{enumerate}[nosep,label=(\roman*)]
  \item If $f'(x) > 0$ for all $x \in (a, b)$, then $f$ is strictly
    increasing on $[a, b]$.
  \item If $f'(x) = 0$ for all $x \in (a, b)$, then $f$ is constant on
    $[a, b]$.
\end{enumerate}
\end{corollary}

\begin{proof}
(i) For $a \le x_1 < x_2 \le b$, the MVT gives $c \in (x_1, x_2)$ with
$f(x_2) - f(x_1) = f'(c)(x_2 - x_1) > 0$.

(ii) For any $x_1, x_2 \in [a, b]$, the MVT gives
$|f(x_2) - f(x_1)| = |f'(c)||x_2 - x_1| = 0$.
\end{proof}

\begin{rlconnection}[title={MVT in optimization convergence analysis}]
The mean value theorem is the workhorse of optimization convergence
proofs. For gradient descent with step size $\eta$ on a function $f$
with $L$-Lipschitz gradient ($\|f'(x) - f'(y)\| \le L\|x-y\|$), the
MVT (generalized to $\R^n$) gives
\[
  f(x - \eta f'(x)) \le f(x) - \eta \|f'(x)\|^2
  + \frac{L\eta^2}{2}\|f'(x)\|^2.
\]
Choosing $\eta = 1/L$ yields $f(x_{k+1}) \le f(x_k) -
\frac{1}{2L}\|f'(x_k)\|^2$, the standard descent lemma that drives
$O(1/k)$ convergence rates for smooth convex functions.
\end{rlconnection}

\begin{example}
\label{ex:mvt_bound}
Let $f(x) = e^x$ on $[0, h]$ for $h > 0$. The MVT gives $c \in (0, h)$
with $e^h - 1 = e^c \cdot h$. Since $1 < e^c < e^h$, we obtain the
bounds $h < e^h - 1 < h\,e^h$.
\end{example}

%% ============================================================
%% SECTION 6: TAYLOR'S THEOREM
%% ============================================================
\section{Taylor's Theorem}
\label{sec:taylor}

\begin{theorem}[Taylor's Theorem with Lagrange Remainder]
\label{thm:taylor}
Let $f \colon [a, b] \to \R$ be such that $f, f', \ldots, f^{(n)}$ are
continuous on $[a, b]$ and $f^{(n+1)}$ exists on $(a, b)$. Then for every
$x \in [a, b]$, there exists $c$ between $a$ and $x$ such that
\[
  f(x) = \sum_{k=0}^{n} \frac{f^{(k)}(a)}{k!}(x - a)^k
  + \frac{f^{(n+1)}(c)}{(n+1)!}(x - a)^{n+1}.
\]
\end{theorem}

\begin{proof}
Define the remainder $R_n(x) = f(x) - \sum_{k=0}^n
\frac{f^{(k)}(a)}{k!}(x-a)^k$. We seek $c$ with
$R_n(x) = \frac{f^{(n+1)}(c)}{(n+1)!}(x-a)^{n+1}$.

Fix $x \neq a$ and define $M$ by $R_n(x) = M(x - a)^{n+1}$, i.e.,
$M = R_n(x)/(x - a)^{n+1}$. Define
\[
  g(t) = f(t) - \sum_{k=0}^{n} \frac{f^{(k)}(a)}{k!}(t - a)^k
  - M(t - a)^{n+1}.
\]
Then $g(a) = 0$ (all terms vanish at $t = a$) and $g(x) = R_n(x) -
M(x - a)^{n+1} = 0$ by the choice of $M$.

By Rolle's theorem, there exists $c_1$ between $a$ and $x$ with
$g'(c_1) = 0$. Now $g'(a) = 0$ (since the first $n$ Taylor terms are
chosen to match derivatives at $a$, and the $(n+1)$-th term contributes
$(n+1)M(t-a)^n$ which vanishes at $a$). Applying Rolle's theorem to
$g'$ on the interval between $a$ and $c_1$, we get $c_2$ with
$g''(c_2) = 0$. Continuing $n+1$ times, we obtain $c = c_{n+1}$ between
$a$ and $x$ with $g^{(n+1)}(c) = 0$. But
$g^{(n+1)}(t) = f^{(n+1)}(t) - (n+1)!\,M$,
so $f^{(n+1)}(c) = (n+1)!\,M$, giving
$R_n(x) = \frac{f^{(n+1)}(c)}{(n+1)!}(x - a)^{n+1}$.
\end{proof}

\begin{example}
\label{ex:taylor_exp}
For $f(x) = e^x$ centered at $a = 0$, $f^{(k)}(0) = 1$ for all $k$,
so the Taylor polynomial is $\sum_{k=0}^n x^k / k!$ and the remainder
is $R_n(x) = e^c x^{n+1}/(n+1)!$ for some $c$ between $0$ and $x$.
For any fixed $x$, $|R_n(x)| \le e^{|x|} |x|^{n+1}/(n+1)! \to 0$
as $n \to \infty$ (by the ratio test applied to $|x|^n/n!$), confirming
$e^x = \sum_{k=0}^\infty x^k/k!$.
\end{example}

\begin{example}
\label{ex:taylor_quadratic}
The first-order Taylor expansion of a smooth function $f$ at $x_0$ is
$f(x) = f(x_0) + f'(x_0)(x - x_0) + \frac{f''(c)}{2}(x - x_0)^2$.
In optimization, setting $x = x_0 - \eta f'(x_0)$ and bounding
$|f''(c)| \le L$ gives
\[
  f(x_0 - \eta f'(x_0)) \le f(x_0) - \eta (f'(x_0))^2
  + \frac{L\eta^2}{2}(f'(x_0))^2,
\]
the descent inequality used to analyze gradient descent.
\end{example}

\begin{rlconnection}[title={Taylor approximations in optimization convergence}]
Taylor's theorem provides the mathematical foundation for convergence
analysis of first- and second-order optimization methods. The
second-order expansion $f(x + h) \approx f(x) + f'(x)h +
\frac{1}{2}f''(x)h^2$ is used to:
\begin{itemize}[nosep]
  \item Derive the \emph{descent lemma} for gradient descent with
    $L$-smooth functions.
  \item Analyze \emph{Newton's method}, which uses the full quadratic
    model to achieve quadratic convergence near a minimizer.
  \item Justify \emph{natural gradient} and \emph{trust region} methods
    in policy optimization, where the policy is locally approximated by
    a quadratic model in the parameter space.
\end{itemize}
\end{rlconnection}

%% ============================================================
%% SECTION 7: RIEMANN INTEGRATION
%% ============================================================
\section{Riemann Integration}
\label{sec:riemann}

\begin{definition}[Partition, Upper and Lower Sums]
\label{def:partition}
Let $f \colon [a, b] \to \R$ be bounded.
\begin{enumerate}[nosep,label=(\roman*)]
  \item A \emph{partition} of $[a, b]$ is a finite set
    $P = \{x_0, x_1, \ldots, x_n\}$ with
    $a = x_0 < x_1 < \cdots < x_n = b$.
  \item On each subinterval $[x_{i-1}, x_i]$, define
    $m_i = \inf\{f(x) : x \in [x_{i-1}, x_i]\}$ and
    $M_i = \sup\{f(x) : x \in [x_{i-1}, x_i]\}$.
  \item The \emph{lower sum} is
    $L(f, P) = \sum_{i=1}^n m_i (x_i - x_{i-1})$.
  \item The \emph{upper sum} is
    $U(f, P) = \sum_{i=1}^n M_i (x_i - x_{i-1})$.
\end{enumerate}
\end{definition}

\begin{definition}[Refinement]
\label{def:refinement}
A partition $Q$ is a \emph{refinement} of $P$ if $P \subseteq Q$.
\end{definition}

\begin{lemma}[Refinement Lemma]
\label{lem:refinement}
If $Q$ is a refinement of $P$, then
$L(f, P) \le L(f, Q) \le U(f, Q) \le U(f, P)$.
\end{lemma}

\begin{proof}
It suffices to prove the result when $Q$ has one additional point $x^*$
in $[x_{j-1}, x_j]$; the general case follows by induction. The lower
sum changes only on $[x_{j-1}, x_j]$:
\begin{align*}
  &\inf_{[x_{j-1}, x^*]} f \cdot (x^* - x_{j-1})
  + \inf_{[x^*, x_j]} f \cdot (x_j - x^*) \\
  &\qquad \ge \inf_{[x_{j-1}, x_j]} f \cdot
  \bigl((x^* - x_{j-1}) + (x_j - x^*)\bigr)
  = m_j(x_j - x_{j-1}),
\end{align*}
since $\inf_{[x_{j-1}, x^*]} f \ge \inf_{[x_{j-1}, x_j]} f$ and
similarly for $[x^*, x_j]$. Hence $L(f, Q) \ge L(f, P)$. The argument
for $U(f, Q) \le U(f, P)$ is analogous (with sup and the inequality
reversed). Finally, $L(f, Q) \le U(f, Q)$ since $m_i \le M_i$ on each
subinterval.
\end{proof}

\begin{corollary}
\label{cor:lower_upper}
For any two partitions $P_1, P_2$ of $[a, b]$:
$L(f, P_1) \le U(f, P_2)$.
\end{corollary}

\begin{proof}
Let $Q = P_1 \cup P_2$ (a common refinement). Then
$L(f, P_1) \le L(f, Q) \le U(f, Q) \le U(f, P_2)$
by Lemma~\ref{lem:refinement}.
\end{proof}

\begin{definition}[Riemann Integral]
\label{def:riemann}
Let $f \colon [a, b] \to \R$ be bounded. Define the \emph{lower
integral} and \emph{upper integral}:
\[
  \underline{\int_a^b} f = \sup_P L(f, P), \qquad
  \overline{\int_a^b} f = \inf_P U(f, P),
\]
where the supremum and infimum are over all partitions $P$ of $[a, b]$.
By Corollary~\ref{cor:lower_upper},
$\underline{\int_a^b} f \le \overline{\int_a^b} f$.
We say $f$ is \emph{Riemann integrable} on $[a, b]$ if
$\underline{\int_a^b} f = \overline{\int_a^b} f$, and the common value
is denoted $\int_a^b f(x)\,dx$ or $\int_a^b f$.
\end{definition}

\begin{theorem}[Integrability Criterion]
\label{thm:int_criterion}
A bounded function $f$ on $[a, b]$ is Riemann integrable if and only if
for every $\varepsilon > 0$ there exists a partition $P$ such that
$U(f, P) - L(f, P) < \varepsilon$.
\end{theorem}

\begin{proof}
$(\Rightarrow)$ Suppose $f$ is integrable with
$\underline{\int} f = \overline{\int} f = I$. Given $\varepsilon > 0$,
by definition of supremum and infimum, choose $P_1$ with
$L(f, P_1) > I - \varepsilon/2$ and $P_2$ with
$U(f, P_2) < I + \varepsilon/2$. Let $P = P_1 \cup P_2$. Then
$L(f, P) \ge L(f, P_1) > I - \varepsilon/2$ and
$U(f, P) \le U(f, P_2) < I + \varepsilon/2$, so
$U(f, P) - L(f, P) < \varepsilon$.

$(\Leftarrow)$ For each $n$, choose $P_n$ with
$U(f, P_n) - L(f, P_n) < 1/n$. Then
$\overline{\int} f - \underline{\int} f \le U(f, P_n) - L(f, P_n) < 1/n$
for all $n$. Hence $\overline{\int} f = \underline{\int} f$.
\end{proof}

\begin{theorem}[Continuous Functions Are Integrable]
\label{thm:cont_integrable}
If $f \colon [a, b] \to \R$ is continuous, then $f$ is Riemann
integrable.
\end{theorem}

\begin{proof}
Since $[a, b]$ is compact and $f$ is continuous, $f$ is uniformly
continuous (Lesson~2, Theorem~7.4). Given $\varepsilon > 0$, choose
$\delta > 0$ such that $|f(x) - f(y)| < \varepsilon/(b - a)$ whenever
$|x - y| < \delta$. Take any partition $P = \{x_0, \ldots, x_n\}$ with
mesh $\|P\| = \max_i(x_i - x_{i-1}) < \delta$. Then on each
$[x_{i-1}, x_i]$, $M_i - m_i < \varepsilon/(b - a)$ (since the
subinterval has diameter less than $\delta$), so
\[
  U(f, P) - L(f, P) = \sum_{i=1}^n (M_i - m_i)(x_i - x_{i-1})
  < \frac{\varepsilon}{b - a} \sum_{i=1}^n (x_i - x_{i-1})
  = \varepsilon.
\]
By Theorem~\ref{thm:int_criterion}, $f$ is integrable.
\end{proof}

\begin{example}
\label{ex:integral_x_squared}
Let $f(x) = x^2$ on $[0, 1]$. Using the uniform partition
$P_n = \{0, 1/n, 2/n, \ldots, 1\}$:
\[
  L(f, P_n) = \sum_{i=1}^n \left(\frac{i-1}{n}\right)^2 \cdot \frac{1}{n}
  = \frac{1}{n^3}\sum_{i=0}^{n-1} i^2
  = \frac{(n-1)n(2n-1)}{6n^3} \to \frac{1}{3}.
\]
Similarly, $U(f, P_n) = \frac{n(n+1)(2n+1)}{6n^3} \to 1/3$. Since
$U(f, P_n) - L(f, P_n) \to 0$, we confirm $\int_0^1 x^2\,dx = 1/3$.
\end{example}

%% ============================================================
%% SECTION 8: PROPERTIES OF THE INTEGRAL
%% ============================================================
\section{Properties of the Integral}
\label{sec:integral_props}

\begin{theorem}[Properties of the Riemann Integral]
\label{thm:integral_props}
Let $f, g$ be Riemann integrable on $[a, b]$ and $\alpha \in \R$. Then:
\begin{enumerate}[nosep,label=(\roman*)]
  \item \textbf{Linearity:}
    $\int_a^b (\alpha f + g) = \alpha \int_a^b f + \int_a^b g$.
  \item \textbf{Monotonicity:} If $f(x) \le g(x)$ for all
    $x \in [a, b]$, then $\int_a^b f \le \int_a^b g$.
  \item \textbf{Additivity:} For $a < c < b$,
    $\int_a^b f = \int_a^c f + \int_c^b f$.
  \item \textbf{Triangle inequality:}
    $\left|\int_a^b f\right| \le \int_a^b |f|$.
\end{enumerate}
\end{theorem}

\begin{proof}
(i) For the scalar multiple: on each subinterval, $\sup(\alpha f) =
\alpha \sup f$ if $\alpha \ge 0$ (and $\alpha \inf f$ if $\alpha < 0$),
so $U(\alpha f, P) = \alpha\, U(f, P)$ for $\alpha \ge 0$. The case
$\alpha < 0$ uses $U(\alpha f, P) = \alpha\, L(f, P)$. Taking
suprema/infima over partitions gives
$\int \alpha f = \alpha \int f$.

For the sum: on each subinterval $[x_{i-1}, x_i]$,
$\inf f + \inf g \le \inf(f + g) \le \sup(f + g) \le \sup f + \sup g$,
so $L(f, P) + L(g, P) \le L(f+g, P)$ and
$U(f+g, P) \le U(f, P) + U(g, P)$. Given $\varepsilon > 0$, choose
$P_1, P_2$ with $U(f, P_1) - L(f, P_1) < \varepsilon/2$ and
$U(g, P_2) - L(g, P_2) < \varepsilon/2$. Let $P = P_1 \cup P_2$. Then
\[
  U(f+g, P) - L(f+g, P) \le \bigl(U(f, P) - L(f, P)\bigr)
  + \bigl(U(g, P) - L(g, P)\bigr) < \varepsilon,
\]
so $f + g$ is integrable. The value equals $\int f + \int g$ by
taking limits of upper sums.

(ii) If $f \le g$, then for any partition $P$,
$U(f, P) \le U(g, P)$, so $\int f = \inf_P U(f, P) \le \inf_P U(g, P)
= \int g$.

(iii) Given $\varepsilon > 0$, choose partitions $P'$ of $[a, c]$ and
$P''$ of $[c, b]$ with $U(f, P') - L(f, P') < \varepsilon/2$ and
$U(f, P'') - L(f, P'') < \varepsilon/2$. Then $P = P' \cup P''$ is a
partition of $[a, b]$ with $L(f, P) = L(f, P') + L(f, P'')$ and
$U(f, P) = U(f, P') + U(f, P'')$. The result follows by taking
$\varepsilon \to 0$.

(iv) Since $-|f(x)| \le f(x) \le |f(x)|$, monotonicity gives
$-\int |f| \le \int f \le \int |f|$.
\end{proof}

%% ============================================================
%% SECTION 9: FUNDAMENTAL THEOREM OF CALCULUS
%% ============================================================
\section{Fundamental Theorem of Calculus}
\label{sec:ftc}

\begin{keyresult}[title={Fundamental Theorem of Calculus}]
\begin{theorem}[FTC Part I]
\label{thm:ftc1}
Let $f \colon [a, b] \to \R$ be Riemann integrable. Define
$F(x) = \int_a^x f(t)\,dt$ for $x \in [a, b]$. Then:
\begin{enumerate}[nosep,label=(\roman*)]
  \item $F$ is continuous on $[a, b]$.
  \item If $f$ is continuous at $c \in (a, b)$, then $F$ is
    differentiable at $c$ and $F'(c) = f(c)$.
\end{enumerate}
\end{theorem}

\begin{theorem}[FTC Part II]
\label{thm:ftc2}
If $f \colon [a, b] \to \R$ is continuous and $G$ is any antiderivative
of $f$ (i.e., $G'(x) = f(x)$ for all $x \in (a, b)$, with $G$
continuous on $[a, b]$), then
\[
  \int_a^b f(x)\,dx = G(b) - G(a).
\]
\end{theorem}
\end{keyresult}

\begin{proof}[Proof of FTC Part I]
(i) Since $f$ is bounded, say $|f(t)| \le M$ for $t \in [a, b]$, we
have for $x, y \in [a, b]$ with $x < y$:
\[
  |F(y) - F(x)| = \left|\int_x^y f(t)\,dt\right|
  \le \int_x^y |f(t)|\,dt \le M(y - x).
\]
Given $\varepsilon > 0$, taking $\delta = \varepsilon/M$ gives
$|F(y) - F(x)| < \varepsilon$ whenever $|y - x| < \delta$. Hence $F$
is (uniformly) continuous.

(ii) Suppose $f$ is continuous at $c$. Given $\varepsilon > 0$, choose
$\delta > 0$ with $|f(t) - f(c)| < \varepsilon$ for $|t - c| < \delta$.
For $0 < |h| < \delta$:
\[
  \frac{F(c + h) - F(c)}{h} - f(c)
  = \frac{1}{h}\int_c^{c+h} f(t)\,dt - f(c)
  = \frac{1}{h}\int_c^{c+h} \bigl(f(t) - f(c)\bigr)\,dt.
\]
For $h > 0$:
$\left|\frac{1}{h}\int_c^{c+h}(f(t) - f(c))\,dt\right|
\le \frac{1}{h}\int_c^{c+h}|f(t) - f(c)|\,dt
< \frac{1}{h} \cdot \varepsilon h = \varepsilon$.
The case $h < 0$ is analogous. Hence $F'(c) = f(c)$.
\end{proof}

\begin{proof}[Proof of FTC Part II]
Let $F(x) = \int_a^x f(t)\,dt$. By Part~I, $F'(x) = f(x) = G'(x)$ for
$x \in (a, b)$. Hence $(F - G)'(x) = 0$ on $(a, b)$, so $F - G$ is
constant on $[a, b]$ by Corollary~\ref{cor:monotone_deriv}(ii). Since
$F(a) = 0$, we get $F(x) = G(x) - G(a)$ for all $x \in [a, b]$.
Evaluating at $x = b$:
$\int_a^b f = F(b) = G(b) - G(a)$.
\end{proof}

\begin{example}
\label{ex:ftc}
To compute $\int_0^1 x^3\,dx$: an antiderivative of $f(x) = x^3$ is
$G(x) = x^4/4$. By FTC Part~II,
$\int_0^1 x^3\,dx = G(1) - G(0) = 1/4 - 0 = 1/4$.
\end{example}

\begin{example}
\label{ex:ftc_derivative}
Let $F(x) = \int_0^x e^{-t^2}\,dt$. By FTC Part~I, $F'(x) = e^{-x^2}$.
While $F$ cannot be expressed in closed form using elementary functions,
the FTC guarantees its differentiability and gives the derivative
explicitly. The function $F$ is closely related to the error function
$\operatorname{erf}(x) = \frac{2}{\sqrt{\pi}} F(x)$, ubiquitous in
probability and statistics.
\end{example}

\begin{warning}[title={Limitations of the Riemann integral}]
The Riemann integral has fundamental limitations:
\begin{enumerate}[nosep,label=(\roman*)]
  \item The indicator function $\indicator_\Q$ of the rationals on $[0,1]$
    is not Riemann integrable: $L(f, P) = 0$ and $U(f, P) = 1$ for every
    partition $P$ (since every subinterval contains both rationals and
    irrationals).
  \item Pointwise limits of Riemann integrable functions need not be
    Riemann integrable, and even when they are, the limit of the
    integrals need not equal the integral of the limit.
  \item The Riemann integral does not interact well with measure theory:
    it cannot integrate with respect to general measures.
\end{enumerate}
These limitations motivate the \emph{Lebesgue integral} (Phase~01,
Measure Theory), which resolves all three issues and provides the
foundation for modern probability theory.
\end{warning}

%% ============================================================
%% SECTION 10: EXERCISES
%% ============================================================
\section{Exercises}
\label{sec:exercises}

\begin{exercise}[Differentiation from the definition]
\textnormal{(\(*\))} \\
Using Definition~\ref{def:derivative}, prove that $f(x) = 1/(x+1)$ is
differentiable on $(-1, \infty)$ and compute $f'(x)$.
\end{exercise}

\begin{exercise}[Chain rule application]
\textnormal{(\(*\))} \\
Let $g(x) = \sin(x^2 + 1)$. Compute $g'(x)$ using the chain rule
(Theorem~\ref{thm:chain_rule}), identifying the inner and outer
functions explicitly.
\end{exercise}

\begin{exercise}[MVT in gradient descent]
\textnormal{(\(**\))} \\
Let $f \colon \R \to \R$ be differentiable with $|f'(x) - f'(y)| \le
L|x - y|$ for all $x, y \in \R$ (i.e., $f'$ is Lipschitz with
constant $L$).
\begin{enumerate}[nosep,label=(\alph*)]
  \item Use Taylor's theorem (Theorem~\ref{thm:taylor} with $n = 1$) to
    show that for all $x, h \in \R$:
    $|f(x + h) - f(x) - f'(x)h| \le \frac{L}{2}h^2$.
  \item Deduce that the gradient descent update $x_{k+1} = x_k -
    \frac{1}{L} f'(x_k)$ satisfies
    $f(x_{k+1}) \le f(x_k) - \frac{1}{2L}(f'(x_k))^2$.
\end{enumerate}
\textit{Hint:} For (a), the Lagrange remainder involves $f''(c)$, and
$|f''(c)| \le L$ follows from the Lipschitz condition on $f'$ via the
MVT.
\end{exercise}

\begin{exercise}[Integrability of monotone functions]
\textnormal{(\(**\))} \\
Prove that every monotone function $f \colon [a, b] \to \R$ is Riemann
integrable.

\textit{Hint:} Assume $f$ is increasing. Use the uniform partition
$P_n$ with $n$ equal subintervals, and show
$U(f, P_n) - L(f, P_n) = \frac{b-a}{n}(f(b) - f(a)) \to 0$.
\end{exercise}

\begin{exercise}[FTC and the Gaussian integral]
\textnormal{(\(**\))} \\
Define $F(x) = \int_0^x e^{-t^2/2}\,dt$.
\begin{enumerate}[nosep,label=(\alph*)]
  \item Use FTC Part~I to show $F$ is differentiable and find $F'(x)$.
  \item Show $F$ is strictly increasing on $[0, \infty)$ and bounded
    above. Conclude that $\lim_{x \to \infty} F(x)$ exists.
  \item Use the bound $e^{-t^2/2} \le e^{-t/2}$ for $t \ge 1$ to show
    $F(x) \le 1 + 2e^{-1/2}$ for all $x \ge 0$.
\end{enumerate}
\end{exercise}

\begin{exercise}[Non-integrability and measure theory motivation]
\textnormal{(\(***\))} \\
Let $f = \indicator_\Q$ on $[0,1]$ (i.e., $f(x) = 1$ if $x \in \Q$,
$f(x) = 0$ otherwise).
\begin{enumerate}[nosep,label=(\alph*)]
  \item Prove that $L(f, P) = 0$ and $U(f, P) = 1$ for every partition
    $P$ of $[0,1]$.
  \item Conclude that $f$ is not Riemann integrable.
  \item The Lebesgue integral assigns $\int_{[0,1]} f\,d\lambda = 0$
    (since $\Q$ has Lebesgue measure zero). Explain informally why this
    is the ``right'' answer and why the Riemann integral fails to
    capture it.
\end{enumerate}
\end{exercise}

%% ============================================================
%% SECTION 11: SUMMARY
%% ============================================================
\section{Summary}
\label{sec:summary}

\begin{itemize}
  \item The \emph{derivative} $f'(c) = \lim_{h \to 0}(f(c+h) - f(c))/h$
    measures the instantaneous rate of change; differentiability at $c$
    implies continuity at $c$, but not conversely (e.g., $|x|$ at $0$).

  \item The \emph{sum, product, quotient, and chain rules} reduce
    differentiation of complex expressions to differentiation of
    simpler components. The chain rule is the mathematical basis of
    backpropagation.

  \item \emph{Fermat's theorem}: at an interior extremum of a
    differentiable function, the derivative vanishes. This leads to
    \emph{Rolle's theorem} and the \emph{mean value theorem}, which
    connects the derivative to the function's total change.

  \item \emph{Taylor's theorem} with Lagrange remainder provides
    polynomial approximations with explicit error bounds:
    $|R_n(x)| \le \frac{M}{(n+1)!}|x - a|^{n+1}$ where
    $M = \sup |f^{(n+1)}|$. This underpins convergence analysis of
    gradient descent and Newton's method.

  \item The \emph{Riemann integral} is defined via upper and lower sums
    over partitions. A bounded function is integrable if and only if for
    every $\varepsilon > 0$ there exists a partition with
    $U(f,P) - L(f,P) < \varepsilon$.

  \item \emph{Continuous functions are Riemann integrable}, via uniform
    continuity on compact sets.

  \item \emph{FTC Part I}: the integral $F(x) = \int_a^x f$ is
    continuous, and $F'(c) = f(c)$ at points of continuity of $f$.
    \emph{FTC Part II}: $\int_a^b f = G(b) - G(a)$ for any
    antiderivative $G$.

  \item The Riemann integral has fundamental limitations (non-integrable
    functions, poor behavior under limits) that motivate the
    \emph{Lebesgue integral} developed in Phase~01.
\end{itemize}

%% ============================================================
%% SECTION 12: FURTHER READING
%% ============================================================
\section{Further Reading}
\label{sec:further_reading}

\begin{itemize}
  \item \textbf{Lesson 4} (Sequences and Series of Functions) extends
    differentiation and integration to function sequences, addressing
    when differentiation and integration commute with limits---a
    question whose definitive answer comes from the Lebesgue theory.

  \item In \textbf{Phase~01} (Measure Theory, Lessons~4--5), the Riemann
    integral is replaced by the Lebesgue integral, which integrates a
    much larger class of functions and satisfies powerful limit theorems
    (monotone convergence, dominated convergence) that the Riemann
    integral lacks.

  \item In \textbf{Phase~01} (Optimization, Lessons~8--10), the mean
    value theorem and Taylor's theorem are extended to $\R^n$ and used
    systematically to analyze gradient descent, proximal methods, and
    interior-point methods.

  \item For multivariable differentiation (partial derivatives, the
    Jacobian, and the general chain rule), see \citet[Chapter 9]{rudin1976}
    or \citet[Chapter 5]{pugh2015}. These are essential for understanding
    backpropagation in neural networks with vector-valued layers.

  \item \citet[Chapter 6]{bartle2011} provides an extensive collection
    of examples and counterexamples for Riemann integrability, clarifying
    the boundary between integrable and non-integrable functions.

  \item For historical context, the development of the integral from
    Cauchy through Riemann to Lebesgue is discussed in
    \citet[Chapter 1]{folland1999}, providing motivation for the modern
    measure-theoretic approach.
\end{itemize}

\bibliography{real-analysis-refs}

\end{document}
